\section*{Hoofdstuk \ref{ch:g6:symmetry} --- Spiegelsymmetrie}

\begin{solution}{exo:g6:symmetry:1}
$A'$ ligt twee vakjes \emph{rechts} van $d$, op dezelfde rij als $A$; $B'$ ligt
vijf vakjes rechts en één vakje hoger --- spiegelposities.
\end{solution}

\begin{solution}{exo:g6:symmetry:2}
De loodlijn op $d$ door $M$, met voetpunt $H$; dan $M'$ aan de andere kant met
$HM' = MH = 3$ cm. De afstand $MM'$ is $3 + 3 = 6$ cm.
\end{solution}

\begin{solution}{exo:g6:symmetry:3}
Verticale as (zoals A): A, H, I, M, O, T, U, V, W, X, Y.
Horizontale as (zoals B): B, C, D, E, H, I, K, O, X.
Beide: H, I, O, X.
(De precieze antwoorden hangen een beetje af van hoe je de letters tekent.)
\end{solution}

\begin{solution}{exo:g6:symmetry:4}
\omterm{def:g6:shapes:triangles}{Gelijkzijdige driehoek}: $3$ assen. Rechthoek: $2$ (de \omterm{def:g6:symmetry:bisector}{middelloodlijnen} van zijn
zijden). \omterm{def:g6:shapes:quadrilaterals}{Ruit}: $2$ (haar diagonalen). Vierkant: $4$. Cirkel: oneindig veel (elke
\omterm{def:g6:lines:objects}{lijn} door het middelpunt).
\end{solution}

\begin{solution}{exo:g6:symmetry:5}
\emph{Niet waar.} Een rechthoek die geen vierkant is, langs een diagonaal
vouwen laat de twee \omterm{def:g3:division:half}{helften} niet samenvallen: de twee driehoeken hebben dezelfde
vorm, maar liggen anders --- probeer het met een blad papier. Alleen bij een
vierkant werken de diagonalen als assen.
\end{solution}

\begin{solution}{exo:g6:symmetry:6}
Spiegelingen bewaren lengtes, hoeken en omtrekken
(\cref{prop:g6:symmetry:preserved}): $A'B' = 4$ cm,
$\widehat{B'A'C'} = 70^\circ$, en de \omterm{def:g4:measure:perimeter}{omtrek} van $A'B'C'$ is $12$ cm.
\end{solution}

\begin{solution}{exo:g6:symmetry:7}
Constructie met de passer zoals in \cref{met:g6:symmetry:bisector}. Voor elk
punt $M$ op de \omterm{def:g6:lines:objects}{lijn} bevestigt de lat dat $MA = MB$
(\cref{prop:g6:symmetry:bisector}).
\end{solution}

\begin{solution}{exo:g6:symmetry:8}
Alle punten op gelijke afstand van $A$ en $B$ vormen de \omterm{def:g6:symmetry:bisector}{middelloodlijn} van
$[AB]$ (\cref{prop:g6:symmetry:bisector}). De fontein mag overal op die \omterm{def:g6:lines:objects}{lijn}
staan (binnen de tuin).
\end{solution}

\begin{solution}{exo:g6:symmetry:9}
Elk \omterm{def:g5:solids:def}{hoekpunt} wordt \omterm{def:g6:lines:perp}{loodrecht} op de as gespiegeld: bij de as van $45^\circ$ landt
een punt $k$ vakjes rechts van de as $k$ vakjes erboven (en omgekeerd) --- de
horizontale en verticale verplaatsingen worden verwisseld. Spiegel de \omterm{def:g5:solids:def}{hoekpunten}
één voor één en verbind ze daarna op volgorde.
\end{solution}

\begin{solution}{exo:g6:symmetry:10}
Het \omterm{def:g6:symmetry:def}{spiegelbeeld} van $M$ is het tweede snijpunt van de cirkel met de loodlijn op
$d$ door $M$ --- het ligt op de cirkel. De reden: de spiegeling bewaart
afstanden en houdt $O$ op zijn plaats (dat ligt op $d$), dus ligt het beeld van
elk punt op afstand $r$ van $O$ weer op afstand $r$ van $O$: de cirkel wordt op
zichzelf afgebeeld.
\end{solution}

\begin{solution}{exo:g6:symmetry:11}
Elk snijpunt $P$ is met dezelfde passeropening $\rho$ vanuit $A$ en vanuit $B$
getekend: dus $PA = \rho = PB$. Volgens
\cref{prop:g6:symmetry:bisector} ligt een punt op gelijke afstand van $A$ en
$B$ op de \omterm{def:g6:symmetry:bisector}{middelloodlijn} van $[AB]$; twee zulke punten bepalen de \omterm{def:g6:lines:objects}{lijn}.
\end{solution}

\begin{solution}{exo:g6:symmetry:12}
Voor elk stuitpunt $P$ op de wand geldt $PB = PB'$ (de spiegeling in $d$ bewaart
afstanden), dus is de padlengte $AP + PB = AP + PB'$. De gebroken \omterm{def:g6:lines:objects}{lijn}
$A \to P \to B'$ is het kortst als hij recht is, dus als $P$ op het \omterm{def:g6:lines:objects}{lijnstuk}
$[AB']$ ligt. Het beste stuitpunt is dus het snijpunt van $[AB']$ met de wand.
\end{solution}

\begin{solution}{pb:g6:symmetry:1}
\textbf{1.} Constructie zoals in \cref{met:g6:symmetry:bisector}; voor elk punt
$M$ op de \omterm{def:g6:lines:objects}{lijn} bevestigt de lat dat $MA = MB$
(\cref{prop:g6:symmetry:bisector}).

\textbf{2.} De fonteinvraag heeft de tweede richting nodig: \emph{alle} punten
op gelijke afstand liggen op de \omterm{def:g6:symmetry:bisector}{middelloodlijn}, dus zijn de mogelijke plekken
precies die \omterm{def:g6:lines:objects}{lijn} en niets anders. En een punt dat niet op de \omterm{def:g6:symmetry:bisector}{middelloodlijn} ligt,
ligt dus \emph{niet} op gelijke afstand: het ligt strikt dichter bij een van de
twee bomen.

\textbf{3.} De twee paden delen het stuk $[MP]$, en $PB = PA$ omdat $P$ op de
\omterm{def:g6:symmetry:bisector}{middelloodlijn} ligt. Het pad $M \to P \to B$ is dus even lang als het pad
$M \to P \to A$, namelijk $MP + PA$. Nu is $MB$ de rechte weg naar $B$ --- die is
hier juist het pad via $P$, dus $MB = MP + PA$; en de rechte weg naar $A$ is
korter dan de omweg via $P$: $MA < MP + PA = MB$. Elk punt aan de kant van $A$
ligt dichter bij $A$.

\textbf{4.} De grens is de \omterm{def:g6:symmetry:bisector}{middelloodlijn} van $[AB]$: school $A$ bedient precies
de \omterm{def:g3:division:half}{helft} van het dorp aan haar kant van die \omterm{def:g6:lines:objects}{lijn} (vraag 3), school $B$ de andere
\omterm{def:g3:division:half}{helft}, en kinderen die op de \omterm{def:g6:lines:objects}{lijn} zelf wonen, mogen kiezen.

\textbf{5.} $O$ ligt op de \omterm{def:g6:symmetry:bisector}{middelloodlijn} van $[AB]$, dus $OA = OB$; en op de
\omterm{def:g6:symmetry:bisector}{middelloodlijn} van $[BC]$, dus $OB = OC$
(\cref{prop:g6:symmetry:bisector}).

\textbf{6.} Uit $OA = OB$ en $OB = OC$ volgt $OA = OC$. Volgens de tweede
richting van \cref{prop:g6:symmetry:bisector} ligt een punt op gelijke afstand
van $A$ en $C$ op de \omterm{def:g6:symmetry:bisector}{middelloodlijn} van $[AC]$: de derde \omterm{def:g6:symmetry:bisector}{middelloodlijn} gaat
ongevraagd door $O$. De drie \omterm{def:g6:symmetry:bisector}{middelloodlijnen} van elke driehoek gaan dus door
\emph{één} punt.

\textbf{7.} De cirkel met middelpunt $O$ en \omterm{def:g3:shapes:shapes}{radius} $OA$ bestaat uit alle punten
op afstand $OA$ van $O$ (\cref{def:g6:lines:circle}); omdat $OB = OA$ en
$OC = OA$, liggen $B$ en $C$ erop. Eén middelpunt, één \omterm{def:g3:shapes:shapes}{radius}, drie punten
bediend: het is de cirkel door $A$, $B$ en $C$.

\textbf{8.} Twee \omterm{def:g6:symmetry:bisector}{middelloodlijnen} zijn genoeg, omdat de derde automatisch volgt
(vraag 6): hun snijpunt voldoet al aan alle drie de gelijkheden van afstanden.
De cirkel teken je met middelpunt $O$ en de passer geopend van $O$ tot een van de
drie \omterm{def:g5:solids:def}{hoekpunten}.

\textbf{9.} Recept: zet drie punten $A$, $B$ en $C$ op de overgebleven boog;
teken de \omterm{def:g6:symmetry:bisector}{middelloodlijnen} van de koorden $[AB]$ en $[BC]$; hun snijpunt is het
middelpunt van het bord, en de afstand tot $A$ is de \omterm{def:g3:shapes:shapes}{radius}. Het werkt omdat de
rand van het bord een cirkel door de drie gemarkeerde punten is --- en vraag 8
vindt \emph{de} cirkel door die punten. De testboog komt precies terug.

\textbf{10.} Liggen $A$, $B$ en $C$ op één \omterm{def:g6:lines:objects}{lijn}, dan staan de \omterm{def:g6:symmetry:bisector}{middelloodlijn} van
$[AB]$ en die van $[BC]$ beide \omterm{def:g6:lines:perp}{loodrecht} op dezelfde \omterm{def:g6:lines:objects}{lijn} $(AB) = (BC)$ --- ze
zijn dus \omterm{def:g6:lines:perp}{parallel} aan elkaar (\cref{prop:g6:lines:facts}) en snijden elkaar, omdat
ze verschillend zijn, nooit. Geen enkel punt ligt op gelijke afstand van alle
drie, en er gaat dus geen cirkel door drie punten op één \omterm{def:g6:lines:objects}{lijn}.

\textbf{11.} Regelmatige \omterm{def:g4:geometry:polygon}{vijfhoek}: $5$ assen, elk door één \omterm{def:g5:solids:def}{hoekpunt} en het midden
van de zijde ertegenover (\omterm{def:g3:numbers:evenodd}{oneven} aantal: elke as is van dat ene soort).
Regelmatige \omterm{def:g4:geometry:polygon}{zeshoek}: $6$ assen --- drie door \omterm{def:g5:solids:def}{hoekpunten} die tegenover elkaar
liggen en drie door de middens van zijden die tegenover elkaar liggen (\omterm{def:g3:numbers:evenodd}{even}
aantal: twee soorten, half en half). Het patroon: een regelmatige \omterm{def:g4:geometry:polygon}{veelhoek} met
$n$ zijden heeft $n$ assen.

\textbf{12.} De knip gaat in één keer door beide lagen, dus verliezen de twee
lagen precies overeenkomstige stukken; openvouwen is de spiegeling in de vouwlijn,
die de rand van de ene laag op die van de andere afbeeldt. Het gat is zijn eigen
\omterm{def:g6:symmetry:def}{spiegelbeeld}: de vouwlijn is zijn \omterm{def:g6:symmetry:axis}{symmetrieas} --- wat je ook geknipt hebt.

\textbf{13.} Het opengevouwen gat heeft minstens $2$ symmetrieassen, de twee
vouwlijnen, en de knip verschijnt in $4$ kopieën (elk openvouwen verdubbelt het
plaatje: $1 \to 2 \to 4$). Vaker vouwen, zoals bij een sneeuwvlok, vermenigvuldigt
de kopieën en de assen verder.

\textbf{14.} $F''$ heeft \emph{dezelfde} omloopsrichting als $F$ (twee
spiegelingen heffen het spiegelen op), dezelfde maat en dezelfde hoogte --- het is
gewoon $F$, naar rechts opgeschoven. De verschuiving is $8$ vakjes: precies
\emph{twee keer} de afstand tussen de twee assen. Twee parallelle spiegels samen
spiegelen dus helemaal niet: ze \emph{schuiven op}.

\textbf{15.} $F''$ heeft weer dezelfde maat, maar staat nu \emph{op zijn kop},
een halve draai gedraaid om het snijpunt van de twee assen: elk punt van $F''$
ligt diametraal tegenover het overeenkomstige punt van $F$ door dat middelpunt.
Twee loodrechte spiegels stellen zich samen tot de \emph{halve draai} --- de
puntspiegeling van \cref{ch:g7:central}.
\end{solution}
